\chapter{Revisão de Literatura}\label{chap:revisao}

Esta seção apresenta os fundamentos teóricos e os principais estudos relacionados ao tema deste trabalho. Inicialmente, descreve-se a arquitetura da Internet das Coisas e suas principais vulnerabilidades de segurança. Em seguida, aborda-se o uso de Aprendizado de Máquina para detecção de intrusão em redes IoT, seguido da caracterização do \textit{dataset} IoT-23 utilizado nos experimentos. Por fim, são apresentados os principais trabalhos relacionados identificados na literatura.

% ══════════════════════════════════════════════════
\section{Internet das Coisas: Arquitetura e Vulnerabilidades}\label{sec:iot}

A Internet das Coisas (\textit{Internet of Things} -- IoT) refere-se à interconexão de dispositivos físicos -- sensores, atuadores, câmeras, eletrodomésticos e equipamentos industriais -- à internet, permitindo a coleta, troca e análise de dados de forma automatizada \citep{Alfahaid2025}. A arquitetura IoT é tipicamente descrita em três camadas principais \citep{Alshahrani2020}:

\begin{itemize}[leftmargin=5em]
    \item \textbf{Camada de percepção:} responsável pela coleta de dados do ambiente físico por meio de sensores e dispositivos embarcados. É a camada mais próxima do hardware e, portanto, mais suscetível a ataques físicos e de \textit{firmware};
    \item \textbf{Camada de rede:} responsável pela transmissão dos dados coletados. Utiliza protocolos como Wi-Fi, ZigBee, Bluetooth, LTE e 5G para comunicação sem fio, e protocolos de aplicação como MQTT (\textit{Message Queuing Telemetry Transport}), CoAP (\textit{Constrained Application Protocol}) e HTTP para troca de mensagens;
    \item \textbf{Camada de aplicação:} responsável pelo processamento e apresentação dos dados ao usuário final, abrangendo plataformas de nuvem, interfaces web e aplicativos móveis.
\end{itemize}

Entre os protocolos utilizados na camada de rede, o MQTT merece destaque por ser amplamente adotado em ambientes IoT devido ao seu baixo consumo de banda e simplicidade. No entanto, uma análise conduzida no mecanismo de busca Shodan revelou até 540.531 vulnerabilidades associadas ao protocolo MQTT em todo o mundo em abril de 2024, evidenciando a dimensão do problema. Por padrão, o MQTT não inclui camada de segurança, sendo responsabilidade do desenvolvedor configurar protocolos como TLS ou SSL \citep{Alshahrani2020}.

As vulnerabilidades em redes IoT decorrem de um conjunto de fatores estruturais. Os dispositivos frequentemente possuem recursos computacionais e de memória reduzidos, o que inviabiliza a execução de mecanismos de segurança complexos. Além disso, muitos equipamentos são comercializados com credenciais padrão que raramente são alteradas pelos usuários finais, e a ausência de atualizações regulares de \textit{firmware} mantém falhas de segurança conhecidas sem correção por longos períodos \citep{Alfahaid2025}.

Os ataques em redes IoT podem ocorrer em diferentes camadas da arquitetura. Na camada de percepção, destacam-se ataques físicos e manipulação de sensores. Na camada de rede, são frequentes ataques de negação de serviço (DoS), ataques intermediário (\textit{man-in-the-middle}) e varreduras de portas. Já na camada de aplicação, ocorrem violações de privacidade e exploração de vulnerabilidades em serviços e aplicações conectadas.

Entre os ataques mais recorrentes em redes IoT, destacam-se:

\begin{itemize}[leftmargin=5em]

    \item \textbf{Botnets IoT:} dispositivos comprometidos são recrutados para compor redes de bots utilizadas em ataques coordenados. O malware Mirai tornou-se um dos exemplos mais conhecidos ao explorar dispositivos IoT com credenciais padrão para execução de ataques DDoS em larga escala;

    \item \textbf{DDoS (\textit{Distributed Denial of Service}):} ataques que utilizam grande volume de tráfego malicioso para tornar serviços indisponíveis;

    \item \textbf{Port Scan:} técnica utilizada para identificar portas abertas e serviços vulneráveis em dispositivos conectados;

    \item \textbf{Comando e Controle (Command and Control -- C\&C):} comunicação estabelecida entre dispositivos infectados e servidores controlados pelos atacantes para coordenação de atividades maliciosas.

\end{itemize}

Diante desse cenário, abordagens tradicionais de segurança baseadas em assinaturas e regras estáticas apresentam limitações importantes em ambientes IoT, principalmente devido à dificuldade de adaptação a ataques inéditos e variantes desconhecidas \citep{Rahman2025}. Nesse contexto, técnicas baseadas em Aprendizado de Máquina vêm sendo investigadas como alternativa para aumentar a capacidade adaptativa dos Sistemas de Detecção de Intrusão em redes IoT.

% ══════════════════════════════════════════════════
\section{Aprendizado de Máquina Aplicado à Segurança em Redes IoT}\label{sec:ml_iot}

O Aprendizado de Máquina (\textit{Machine Learning} -- ML) é um subcampo da inteligência artificial que desenvolve algoritmos capazes de aprender padrões a partir de dados e realizar predições ou classificações sem serem explicitamente programados para cada caso \citep{Alfahaid2025}. No contexto da segurança em redes IoT, o ML é aplicado principalmente para construir Sistemas de Detecção de Intrusão (\textit{Intrusion Detection Systems} -- IDS) baseados em anomalias, que identificam comportamentos suspeitos no tráfego de rede sem depender de assinaturas predefinidas.

As abordagens de ML para IDS em IoT são classificadas em duas grandes categorias:

\begin{itemize}[leftmargin=5em]
    \item \textbf{Aprendizado supervisionado:} o modelo é treinado com dados rotulados, onde cada amostra de tráfego possui uma classe associada (benigno ou tipo de ataque específico). Algoritmos como \textit{Random Forest}, SVM (\textit{Support Vector Machine}), Árvore de Decisão e Redes Neurais Artificiais são amplamente utilizados nessa categoria. A principal vantagem é a alta acurácia quando há dados rotulados suficientes; a limitação é a dependência de um processo de rotulação custoso e a incapacidade de detectar ataques completamente novos;
    \item \textbf{Aprendizado não supervisionado:} o modelo aprende o comportamento normal da rede e sinaliza desvios como potenciais ataques, sem necessidade de dados rotulados. Técnicas como \textit{K-Means}, \textit{Isolation Forest} e \textit{Autoencoders} se enquadram nessa categoria. São mais flexíveis para detectar ameaças desconhecidas, porém tendem a gerar mais falsos positivos.
\end{itemize}

Entre os algoritmos supervisionados mais utilizados na literatura para detecção de ataques IoT, destacam-se:

\begin{itemize}[leftmargin=5em]
    \item \textbf{Random Forest:} conjunto de árvores de decisão treinadas com subconjuntos aleatórios dos dados. É robusto ao sobreajuste (\textit{overfitting}) e lida bem com dados desbalanceados. Experimentos comparativos mostram que o \textit{Random Forest} frequentemente alcança os maiores índices de acurácia em cenários de detecção de botnet em dispositivos IoT.
    \item \textbf{SVM (\textit{Support Vector Machine}):} constrói um hiperplano de separação entre classes no espaço de \textit{features}. Apresenta bom desempenho em problemas de alta dimensionalidade, porém com maior custo computacional em datasets grandes;
    \item \textbf{Árvore de Decisão:} modelo de classificação baseado em regras hierárquicas de decisão. Altamente interpretável e de baixo custo computacional, sendo útil para análise do processo de detecção;
    \item \textbf{CNN (\textit{Convolutional Neural Network}):} rede neural profunda originalmente desenvolvida para processamento de imagens, mas adaptada para classificação de tráfego de rede ao tratar as \textit{features} como sequências. Abordagens híbridas combinando CNN e LSTM têm alcançado até 99\% de acurácia em cenários de detecção de intrusão em redes IoT.
\end{itemize}

Um aspecto crítico a considerar é que a maioria dos estudos de ML para IDS foi desenvolvida e avaliada em datasets de redes convencionais, como NSL-KDD e UNSW-NB15. O tráfego IoT, no entanto, possui características distintas: baixo volume de dados por dispositivo, uso de protocolos específicos e padrões de comunicação periódicos e previsíveis. Essa diferença levanta uma questão central para o presente trabalho: se os algoritmos desenvolvidos para redes convencionais mantêm desempenho satisfatório quando aplicados ao contexto IoT \citep{SouzaArima2024}.

% ══════════════════════════════════════════════════
\section{Dataset IoT-23} \label{sec:dataset}

O IoT-23 é o conjunto de dados adotado neste trabalho. Trata-se
de um dataset de larga escala e disponível publicamente, contendo tráfego de
rede rotulado de dispositivos IoT reais, capturado no Stratosphere Laboratory
da Universidade Técnica Tcheca em Praga. É o primeiro dataset a combinar
execução real de malware em dispositivos IoT físicos com tráfego benigno,
tornando-o representativo de cenários reais de ameaças.

O IoT-23 contém 20 capturas de tráfego malicioso executado em dispositivos
IoT e 3 capturas de tráfego benigno, totalizando 23 cenários. Foi publicado
em janeiro de 2020, com capturas realizadas entre 2018 e 2019 no Stratosphere
Laboratory. Seu objetivo é oferecer um dataset de infecções reais de malware
em IoT e tráfego benigno para que pesquisadores desenvolvam algoritmos de
aprendizado de máquina.

Os cenários de ataque envolvem a execução de malwares como Mirai, Torii,
Hide and Seek, Hajime e outros em dispositivos como Raspberry Pi. O dataset
fornece tráfego de rede rotulado, com características detalhadas e rotulagem
baseada em análise humana, disponibilizando tanto arquivos \texttt{.pcap}
quanto logs de conexão para análise.

Os três tipos de tráfego malicioso mais comuns no dataset são:
PartOfAHorizontalPortScan (213.852.924 fluxos), Okiru (47.381.241 fluxos)
e DDoS (19.538.713 fluxos). Essa distribuição desbalanceada é uma
característica relevante a ser tratada durante o pré-processamento dos
dados neste trabalho.

Embora o IoT-23 permaneça amplamente utilizado como benchmark na
literatura, datasets mais recentes têm sido propostos para cobrir um
espectro maior de ataques. O CIC-IoT2023, desenvolvido pelo Canadian
Institute for Cybersecurity, contempla 33 tipos de ataque organizados em
6 categorias, capturados com dispositivos IoT reais como agentes atacantes
\citep{Neto2023CICIoT}. A adoção do IoT-23 neste trabalho justifica-se
pela existência de baseline direto na literatura nacional — especialmente
\citep{SouzaArima2024} — o que permite comparação direta dos resultados
obtidos.

O dataset está disponível gratuitamente para download em:
\url{https://www.stratosphereips.org/datasets-iot23}
\citep{IoT23Dataset2020}.

% ══════════════════════════════════════════════════
\section{Trabalhos Relacionados}\label{sec:relacionados}

Esta seção apresenta os principais trabalhos que investigaram a detecção de ataques em redes IoT com uso de ML, com ênfase em publicações nacionais recentes, especialmente aquelas apresentadas no Simpósio Brasileiro de Segurança da Informação e de Sistemas Computacionais (SBSeg).

\cite{SouzaArima2024} avaliaram algoritmos de ML para detecção de malware IoT utilizando o mesmo dataset IoT-23 adotado neste trabalho. Os autores compararam \textit{Random Forest}, SVM, Árvore de Decisão e CNN, com o \textit{Random Forest} alcançando a maior acurácia geral e a CNN obtendo boas métricas de precisão e F1-\textit{Score}. Este trabalho constitui o principal \textit{baseline} para os experimentos aqui propostos, sendo o ponto de comparação direto dos resultados a serem obtidos.

\cite{Bochie2020} realizaram uma avaliação ampla de técnicas de ML e aprendizado profundo para detecção de ataques em redes IoT, utilizando \textit{traces} públicos de tráfego. Os autores investigaram modelos como \textit{Random Forest}, SVM, MLP e \textit{Autoencoders}, concluindo que abordagens de aprendizado profundo podem superar os algoritmos clássicos em cenários com maior variedade de ataques. O trabalho é relevante para contextualizar as limitações das abordagens supervisionadas tradicionais.

\cite{Alfahaid2025} publicaram um survey abrangente sobre soluções de ML para segurança em redes IoT, cobrindo publicações de 2020 a 2024. O estudo analisa técnicas como ML, aprendizado profundo, \textit{ensemble learning}, \textit{transfer learning} e \textit{federated learning} aplicadas a problemas de segurança IoT, incluindo detecção de intrusão, proteção de privacidade e gerenciamento de confiança em diferentes domínios como saúde, cidades inteligentes e veículos conectados. Este survey fornece a base teórica mais atualizada para o presente trabalho.

\cite{Rahman2025} apresentaram um survey sobre sistemas de detecção de intrusão em redes IoT, discutindo as principais abordagens, datasets utilizados e desafios em aberto na área. O trabalho é útil para contextualizar o presente trabalho no estado da arte internacional e identificar lacunas de pesquisa ainda não exploradas.

A \cref{tab:relacionados} sintetiza os principais trabalhos relacionados, destacando os algoritmos utilizados, os datasets e os resultados reportados.

A partir da análise dos trabalhos relacionados, observa-se que, embora existam estudos comparativos de algoritmos de ML em redes IoT, ainda há espaço para investigar com maior profundidade quais modelos apresentam melhor equilíbrio entre acurácia, tempo de inferência e robustez a classes desbalanceadas no contexto específico do dataset IoT-23 -- lacuna que este trabalho se propõe a preencher.

\begin{table}[ht]
\centering
\caption{Síntese dos principais trabalhos relacionados.}
\label{tab:relacionados}
\begin{tabular}{p{3.5cm}p{3.5cm}p{2.5cm}p{4cm}}
\hline
\textbf{Trabalho} & \textbf{Algoritmos} & \textbf{Dataset} & \textbf{Principal resultado} \\
\hline
\cite{SouzaArima2024} & RF, SVM, DT, CNN & IoT-23 & RF maior acurácia geral \\
\cite{Bochie2020} & RF, SVM, MLP, Autoencoder & Traces IoT públicos & DL supera ML clássico em ataques variados \\
\cite{Alfahaid2025} & Survey (ML, DL, FL) & Múltiplos & Revisão de tendências 2020--2024 \\
\cite{Rahman2025} & Survey (IDS em IoT) & Múltiplos & Desafios e lacunas em aberto \\
\hline
\end{tabular}
\end{table}